\documentclass[11pt]{article}
\usepackage[utf8]{inputenc}
\usepackage{amsmath,amssymb}
\usepackage{hyperref}
\title{Robust Cell Signaling Pathway Network under Distribution Shift}
\author{Anonymous}
\date{}
\begin{document}
\maketitle

\begin{abstract}
This paper studies Cell Signaling Pathway Network. Grounded in a real, citable literature \cite{ref1,ref2,ref3,ref4,ref5,ref6,ref7,ref8},
we propose a focused method and report \textbf{illustrative} experiments.
All references below are real and verifiable (OpenAlex/arXiv); the reported
numbers are simulated to illustrate intended behaviour.
\end{abstract}

\section{Introduction}
Cell Signaling Pathway Network is an active research area. This work targets a concrete gap identified from
the recent literature and contributes a method together with an evaluation protocol.

\section{Related Work (real literature)}
The following works, retrieved from public bibliographic APIs, frame our study:
\begin{itemize}
\item Jin et al. (2021) — "Inference and analysis of cell-cell communication using CellChat", Nature Communications (cited 8535).
\item Massagué (1998) — "TGF-β SIGNAL TRANSDUCTION", Annual Review of Biochemistry (cited 7574).
\item Jiang et al. (2021) — "Ferroptosis: mechanisms, biology and role in disease", Nature Reviews Molecular Cell Biology (cited 7408).
\item Yarden et al. (2001) — "Untangling the ErbB signalling network", Nature Reviews Molecular Cell Biology (cited 6990).
\item Sica et al. (2012) — "Macrophage plasticity and polarization: in vivo veritas", Journal of Clinical Investigation (cited 6368).
\item Fabregat et al. (2015) — "The Reactome pathway Knowledgebase", Nucleic Acids Research (cited 6019).
\end{itemize}

\section{Method}
We model distribution shift explicitly and optimise a worst-case objective over an uncertainty set of plausible test distributions. We instantiate this idea for Cell Signaling Pathway Network and analyse its cost/quality trade-off.

\section{Experiments (Illustrative)}
\begin{center}
\begin{tabular}{lcc}
System & Primary Metric & Efficiency \\ \hline
Strong baseline & 0.812 & 1.00$\times$ \\
Ablation & 0.828 & 0.92$\times$ \\
\textbf{Ours} & \textbf{0.851} & \textbf{0.71$\times$}
\end{tabular}
\end{center}
\emph{Metrics simulated to illustrate the intended effect; not measured.}

\section{Conclusion}
We connected a real literature to a concrete method for Cell Signaling Pathway Network. Future work will
validate the illustrative results with real experiments.

\begin{thebibliography}{9}
\bibitem{ref1} Suoqin Jin, Christian F. Guerrero‐Juarez, Lihua Zhang et al.. Inference and analysis of cell-cell communication using CellChat. \emph{Nature Communications}, 2021. DOI:10.1038/s41467-021-21246-9
\bibitem{ref2} Joan Massagué. TGF-β SIGNAL TRANSDUCTION. \emph{Annual Review of Biochemistry}, 1998. DOI:10.1146/annurev.biochem.67.1.753
\bibitem{ref3} Xuejun Jiang, Brent R. Stockwell, Marcus Conrad. Ferroptosis: mechanisms, biology and role in disease. \emph{Nature Reviews Molecular Cell Biology}, 2021. DOI:10.1038/s41580-020-00324-8
\bibitem{ref4} Yosef Yarden, Mark X. Sliwkowski. Untangling the ErbB signalling network. \emph{Nature Reviews Molecular Cell Biology}, 2001. DOI:10.1038/35052073
\bibitem{ref5} Antonio Sica, Alberto Mantovani. Macrophage plasticity and polarization: in vivo veritas. \emph{Journal of Clinical Investigation}, 2012. DOI:10.1172/jci59643
\bibitem{ref6} Antonio Fabregat, Konstantinos Sidiropoulos, Phani Garapati et al.. The Reactome pathway Knowledgebase. \emph{Nucleic Acids Research}, 2015. DOI:10.1093/nar/gkv1351
\bibitem{ref7} Alan R. Saltiel, C. Ronald Kahn. Insulin signalling and the regulation of glucose and lipid metabolism. \emph{Nature}, 2001. DOI:10.1038/414799a
\bibitem{ref8} Minoru Kanehisa, Miho Furumichi, Yoko Sato et al.. KEGG: integrating viruses and cellular organisms. \emph{Nucleic Acids Research}, 2020. DOI:10.1093/nar/gkaa970
\end{thebibliography}
\end{document}
